% Chapter 14 worked example. Included by ch14_resources.tex.
\section{Worked Example: Imaging Platform Refresh: Maintain, Refurbish, Replace, or Share}
\label{sec:comp-ch14-worked-example}

\subsection{Purpose and Status}
This worked example demonstrates how the methods in Chapter~14 can be
translated into a reviewable institutional decision. All numerical inputs are
synthetic unless a source is identified explicitly. They are intended to show the
logic of the analysis, not to report empirical findings or establish a universal
threshold. Local use requires replacement of the assumptions with current,
versioned evidence and approval through the relevant clinical, managerial,
economic, legal, technical, and governance processes.

A hospital imaging platform depends on ageing local servers, proprietary interfaces, and a supported but approaching end-of-life software stack. The hospital must compare continued maintenance, modular refurbishment, full replacement, and a shared managed service. 

\subsection{Decision Problem and Comparators}
\textbf{Decision problem.} Which asset strategy preserves safe service, value, resilience, and exit capability over the next five years?

The comparator set is deliberately broader than the existing pathway. This prevents
a new technology from receiving a lower evidentiary threshold than staffing,
workflow, shared-service, conventional digital, or no-change alternatives.

\begin{longtable}{@{}p{0.08\textwidth}p{0.84\textwidth}@{}}
\caption{Comparators for the Chapter 14 worked example}
\label{tab:comp-ch14-comparators}\\
\toprule
\textbf{Option} & \textbf{Comparator or strategy} \\
\midrule
\endfirsthead
\toprule
\textbf{Option} & \textbf{Comparator or strategy} \\
\midrule
\endhead
\bottomrule
\endlastfoot
1 & Maintain current platform for a limited period \\
2 & Modular refurbishment and software modernization \\
3 & Full local replacement \\
4 & Regional shared managed service \\
\end{longtable}

\subsection{Model and Decision Structure}
The analytical structure should follow the decision rather than the available
software. The team should map the causal pathway from intervention input to
information, human decision, action, outcome, resource consequence, and review.
The structure should also identify capacity constraints, uncertainty, subgroup
variation, implementation requirements, and events that would change the
recommendation.

A compact quantitative relationship used in this example is:
\begin{equation}
\mathrm{EAC}_j=\mathrm{PV}(C_j)\left[\frac{r(1+r)^{n_j}}{(1+r)^{n_j}-1}\right]
\label{eq:comp-ch14-key-relationship}
\end{equation}
The equation is an organizing aid rather than a substitute for disaggregated evidence,
qualitative findings, or mandatory safeguards. Its variables and interpretation must
be defined in the local protocol.

\subsection{Synthetic Parameters and Evidence Sources}
Table~\ref{tab:comp-ch14-parameters} provides the base-case inputs. A
production analysis should add parameter distributions, correlations, structural
alternatives, data dates, system versions, and evidence-confidence ratings.

\begin{longtable}{@{}p{0.32\textwidth}p{0.22\textwidth}p{0.38\textwidth}@{}}
\caption{Synthetic parameters for the Chapter 14 worked example}
\label{tab:comp-ch14-parameters}\\
\toprule
\textbf{Parameter} & \textbf{Base value} & \textbf{Source or interpretation} \\
\midrule
\endfirsthead
\toprule
\textbf{Parameter} & \textbf{Base value} & \textbf{Source or interpretation} \\
\midrule
\endhead
\bottomrule
\endlastfoot
Current asset age & 7 years & Asset register \\
Remaining vendor support & 24 months & Support contract \\
Unplanned downtime & 38 hours/year & Service records \\
Five-year PV cost: maintain & EUR 2.4 million & Option model \\
Five-year PV cost: refurbish & EUR 3.1 million & Option model \\
Five-year PV cost: replace & EUR 4.6 million & Option model \\
Five-year PV cost: shared service & EUR 3.8 million & Option model \\
Estimated recoverable residual value & EUR 280,000 & Market and take-back estimate \\
\end{longtable}

\subsection{Step-by-Step Analysis}
The sequence below is intended to make the analysis reproducible and to ensure
that evidence, economics, governance, and implementation develop together.

\begin{longtable}{@{}p{0.07\textwidth}p{0.55\textwidth}p{0.30\textwidth}@{}}
\caption{Analytical sequence}
\label{tab:comp-ch14-analysis-steps}\\
\toprule
\textbf{Step} & \textbf{Required work} & \textbf{Decision record} \\
\midrule
\endfirsthead
\toprule
\textbf{Step} & \textbf{Required work} & \textbf{Decision record} \\
\midrule
\endhead
\bottomrule
\endlastfoot
1 & Define functional service requirements and non-negotiable safety, support, cybersecurity, and continuity conditions. & Evidence and responsible owner to be recorded \\
2 & Compare lifecycle cost and equivalent annual cost only among eligible options. & Evidence and responsible owner to be recorded \\
3 & Assess compatibility, validation burden, downtime, maintainability, software support, energy, residual value, and exit. & Evidence and responsible owner to be recorded \\
4 & Examine reuse, refurbishment, transfer, take-back, secure erasure, and chain of custody. & Evidence and responsible owner to be recorded \\
5 & Test downside scenarios for failure, vendor exit, delayed migration, and demand growth. & Evidence and responsible owner to be recorded \\
6 & Document the transition, data closure, receiving party, final destination, and organizational learning. & Evidence and responsible owner to be recorded \\
\end{longtable}

\subsection{Illustrative Outputs}
The outputs in Table~\ref{tab:comp-ch14-outputs} show how technical,
clinical, operational, economic, equity, and governance findings should be kept
visible rather than compressed into a single favourable or unfavourable score.

\begin{longtable}{@{}p{0.30\textwidth}p{0.24\textwidth}p{0.38\textwidth}@{}}
\caption{Illustrative model and decision outputs}
\label{tab:comp-ch14-outputs}\\
\toprule
\textbf{Output} & \textbf{Result} & \textbf{Interpretation} \\
\midrule
\endfirsthead
\toprule
\textbf{Output} & \textbf{Result} & \textbf{Interpretation} \\
\midrule
\endhead
\bottomrule
\endlastfoot
Maintain option & Lowest apparent cost & Fails supportability beyond 24 months without a transition \\
Refurbish option & Strong short-term value & Depends on software and security support \\
Replace option & Highest initial cost & Best control but substantial embodied resource and transition burden \\
Shared-service option & Moderate cost and resilience potential & Data, dependency, and exit terms are decision critical \\
Preferred strategy & Modular refurbishment with prepared shared-service option & Preserves flexibility \\
\end{longtable}

\subsection{Sensitivity, Uncertainty, and Subgroups}
At minimum, the completed analysis should vary the parameters most likely to
change the decision, test at least one credible alternative model structure, and
report subgroup results separately where performance, access, response, burden,
or opportunity cost may differ. Deterministic analysis should identify thresholds;
probabilistic analysis should be used where credible parameter distributions and
correlations are available. Scenario analysis is preferable when structural or
future uncertainty cannot be represented defensibly by one probability model.

The record should distinguish uncertainty that can be reduced through evidence,
uncertainty that can be managed through safeguards, and uncertainty that remains
too consequential to accept. Missing evidence should not be represented as an
average or neutral value without justification.

\subsection{Interpretation and Recommendation}
Undertake modular refurbishment to preserve safe useful life while preparing a contractually robust shared-service option. Do not extend the current platform beyond supported and validated operation. Use recovery value and take-back only where chain of custody and final disposition are verified.

The recommendation is conditional rather than permanent. It applies only to the
specified population, setting, intervention configuration, evidence cut-off, and
version. The following conditions should be incorporated into the approval,
contract, implementation, monitoring, or renewal record:

\begin{itemize}
 \item support, cybersecurity, and validation remain acceptable throughout the planned period.
 \item equivalent annual cost uses a common boundary and realistic residual value.
 \item data migration, secure erasure, and fallback are tested before transition.
 \item spare parts, documentation, and specialist capability are available.
 \item renewal and replacement decisions are reviewed before the support deadline.
\end{itemize}

\subsection{Decision Statement}
\begin{quote}
For the specified decision problem and comparator set, the institution should
\emph{[proceed / proceed with conditions / redesign / pilot / restrict / defer /
replace / retire]}. The decision applies to \emph{[population, setting, version,
and evidence period]}, is owned by \emph{[accountable authority]}, and will be
reviewed on \emph{[date]} or earlier if \emph{[trigger events]} occur. The
evidence, assumptions, unresolved uncertainty, mandatory safeguards, monitoring
thresholds, and exit arrangements are recorded in the accompanying dossier.
\end{quote}
